\section{Related Works}

\paragraph{LLM Multi-Agent Systems.}
An emerging research focus examines using LLMs~\citep{achiam2023gpt, wang2024comprehensive} as central controllers to develop LLM agents that interact with the external world beyond text domains~\citep{xie2024osworld, deng2024mind2web,zhangoffline,zhang2025nemotron}. 
While single-agent systems~\citep{yao2022react,zhang2024ecoact,zhang2023ecoassistant} excel in specific tasks, they struggle with challenges requiring collaboration and collective intelligence. To address this, studies have explored LLM-powered multi-agent systems, where multiple interactive agents work concurrently~\citep{li2023camel,hong2023metagpt}. 
These systems leverage the specialized skills and roles of individual agents, enabling collaborative problem-solving for complex tasks by simulating real-world cooperation patterns.

\paragraph{LLM for Judging.}
Numerous studies have explored the use of large language models (LLMs) as evaluators to assess various tasks based on pre-defined standards~\citep{gu2024survey,fu2023gptscore,liu2023g,thakur2024judging,li2023alpacaeval,hu2024language}. 
For instance, \citet{zheng2023judging,chan2023chateval} utilize LLMs to evaluate the performance of LLMs in chat conversation scenarios, which would otherwise incur significant labor costs if performed by humans. 
Another notable example is \citet{van2024field,miao2023selfcheck}, who employ LLMs as evaluators in the context of text summarization which also heavily relies on human efforts.
In the field of agentic systems, related research includes \citet{shinn2024reflexion}, who adopt the concept of LLMs-as-judges to analyze task feedback signals and guide corrective actions.
Similarly, \citet{zhuge2024agent} demonstrate the use of LLMs to provide detailed evaluations of agentic systems within their proposed DevAI dataset. Despite these advancements, failure attribution remains a manual process, with evaluation results serving only as a reference for such attributions

\paragraph{Reward Models}
Most reward models (RMs) are designed either to predict human preference rankings for outputs generated by large language models~\citep{zhong2025comprehensive} or to evaluate the reasoning process step by step, rather than assessing only the final answer~\citep{wang2023math,zheng2024processbench,lightman2023let,cui2025process}. A number of studies have proposed training process-level reward models that evaluate the correctness of intermediate reasoning steps produced by a single LLM~\citep{lightman2023let, cui2025process}. For instance, Math-Shepherd~\citep{wang2023math} employs automatically generated supervision data to assign reward scores to each step in solving mathematical problems. Similarly, ProcessBench introduces a benchmark of step-by-step solutions annotated by human experts, identifying the location of errors within mathematical problem-solving processes. In this setting, models are tasked with detecting the earliest erroneous step or confirming that the entire solution is correct. However, these works focus primarily on constructing reward models for evaluating the outputs of individual LLMs, rather than identifying the errors in complex agentic systems.  


\section{Conclusion}

In this study, we propose and formulate a new research area: automated failure attribution in LLM multi-agent systems, an area that has been largely overlooked in current research. 
To advance this field, we introduce the Who\&When dataset, which consists of 127 multi-agent systems with extensive failure logs meticulously annotated with failure details.
Furthermore, we develop and evaluate three automated failure attribution methods, highlighting the challenges and complexities of this task. Our findings underscore the significant difficulty of automated failure attribution and emphasize the urgent need for further research in this emerging area. 